%%%PAGE 14 rot=0 legibility=fair summary=静电场：示波器电子偏转、等效重力场（等效最高点、g等）、静电平衡与静电屏蔽、均匀球壳、静电应用；磁场：地磁场方向
%%%BLOCK id=014-01 ch=09 sec=示波管·电子偏转 kind=knowledge alt=none cont=none y=0.065-0.115
示波器：电子偏转\quad \fbox{$\Delta y\propto U$} % 手写红笔把 Δy∝U 圈起；标题下有波浪线
%%%END
%%%BLOCK id=014-02 ch=09 sec=等效重力场 kind=method alt=04,06 cont=none y=0.10-0.29
\textbf{等效场}

%FIG p014_a 0.01 0.104 0.277 0.29
\notefig[0.3]{figures/p014_a.png}
% 图上标注（其中最高点右侧的 E_p? max 在图外右侧，已被裁掉，故在下行转录）：
% 高 / E_k min / 反重点 / E_p? max ；最低 / E_k max / E_p? min / 同重点 / F' / mg
\textbf{图上标注}：高\ $E_{k}\,\mathrm{min}$\ 反重点\ $E_{p\mbox{\unclear{等}}}\,\mathrm{max}$；\unclear{最}低\ $E_{k}\,\mathrm{max}$\ $E_{p\mbox{\unclear{等}}}\,\mathrm{min}$\ $F'$\ 同重点\ $mg$

完整圆周：等效最高点\quad $v_{\text{高}}\ge\sqrt{g'R}$ % 手写在 V高≥√(g'R) 下有红色波浪线着重

\[
g_{\text{等}}=\frac{F_{\text{等}}}{m}=\frac{\sqrt{(mg)^2+(qE)^2}}{m}
\] % 手写在根号下方有红色波浪线着重
%%%END
%%%BLOCK id=014-03 ch=09 sec=静电平衡 kind=knowledge alt=none cont=none y=0.315-0.44
\textbf{静电平衡}

%FIG p014_b 0.045 0.318 0.21 0.435
\notefig[0.25]{figures/p014_b.png}
% 图中：由“稳定”引出一条弧线箭头，指向标题“静电平衡”
$E_{\text{内}}=0$

$\vec{E}_{\text{感应电荷}}+\vec{E}_{\text{原}}=0$\qquad 导体为等势体，电场线与表面垂直
%%%END
%%%BLOCK id=014-04 ch=09 sec=静电屏蔽 kind=knowledge alt=none cont=none y=0.44-0.67
\textbf{静电屏蔽}
\begin{itemize}
\item 金属导体\quad 屏蔽外来电场线$'$
\item 金属导体 $+$ 壳不接地\quad 无法屏蔽内线$'$
\item 金属导体 $+$ 外壳接地\quad 屏蔽内线
\end{itemize}
%FIG p014_c 0.075 0.486 0.495 0.665
\notefig[0.7]{figures/p014_c.png}
%%%END
%%%BLOCK id=014-05 ch=09 sec=场强·非质点电荷 kind=knowledge alt=none cont=none y=0.50-0.55
均匀球壳$Q$可等效于集中在球心处点电荷\unclear{上白项}
%%%END
%%%BLOCK id=014-06 ch=09 sec=静电的应用 kind=knowledge alt=none cont=none y=0.67-0.80
\textbf{应用}：
\begin{itemize}
\item 静电吸附：吸尘、除尘、复印
\item 尖端\unclear{放}电 % CHECK: 手写字形近“防”，按物理术语“尖端放电”转录
\begin{itemize}
\item 利用：避雷针\ 打火机
\item 避免：电晕、油罐、车接地铁链\ 导电橡胶轮胎
% 此行“油罐、”与“车接地铁链”之间有一块白色修正贴（上有“、”）；其上方另有小字批注及一条竖线指向此行，字迹难辨：
（上方小字批注：\unclear{车软}）
\end{itemize}
\item 静电屏蔽：网状高压工作服，密闭金属罩
\end{itemize}
%%%END
%%%BLOCK id=014-07 ch=11 sec=地磁场 kind=knowledge alt=none cont=none y=0.82-0.96
\textbf{磁场}

%FIG p014_d 0.11 0.828 0.285 0.958
\notefig[0.25]{figures/p014_d.png}
南极\ $B$向上

北极\ $B$向下

赤道\ 平行于地面向北

南半球：$B_1$向北、$B_2$向上

北半球：$B_1$向北、$B_2$向下
%%%END
%%%PAGEEND 14
